% !TeX root = ../../../main.tex
\section{The obstruction in the valuation tower}
\label{sec:as_induction_valuation_tower}

{\color{blue}
\noindent\textbf{Second obstruction: comparability does not suffice.}
To isolate the second difficulty, we temporarily grant that the two
exceptional valuations are compatible and that both proposed inputs
(IH$_{r-1}$) and (AS) hold. We ask whether these assumptions alone suffice
to deduce unramifiedness of the original symmetric-power torsor.
\par}

We may reduce to the case in which $\cP\to U$ is connected. Choose a
connected component $\cP_0\subseteq\cP$, and let $G_0\subseteq G$ be the
subgroup preserving it. Then $\cP_0\to U$ is a connected $G_0$-torsor.
Replacing $(\cP,G,H)$ by $(\cP_0,G_0,H\cap G_0)$ restricts the four quotient
schemes to the corresponding connected components, and both unramifiedness
hypotheses remain valid. We may therefore assume that $\cP$ is connected.
We may also assume that $H\ne0$: if $H=0$, then $\cP/H=\cP$ and $G/H=G$,
so (IH$_{r-1}$) is already the desired conclusion.

Under these reductions the four schemes in the Cartesian square of
\Cref{prop:torsor_tower_diagram} are integral, so we may pass to function
fields. Put
\[
K_C:=\Frac V_C=K(C^{(d)}),\qquad
K_{C'}:=\Frac V_{C'}=K(C'^{(d)}),
\]
and
\[
F:=K((\cP/H)^{[d]}),\qquad
E:=K(\cP^{[d]_G}).
\]
Then $F/K_C$ is Galois with group $G/H$, $E/F$ is Galois with group $H$,
and $E/K_C$ is Galois with group $G$. The Cartesian square identifies the
function field of $\cP^{[d]_H}$ with $E\cdot K_{C'}$. Its generic fibre over
$(\cP/H)^{(d)}$ is the finite domain $E\otimes_FK_{C'}$, hence a field.

We use here the following standard comparison between a Galois extension and
its base change.

\begin{lemma}\label{lem:galois-base-change}
Let $E/F$ be a finite Galois extension with group $H$, let $K/F$ be a field
extension, and form all fields inside a common overfield. Put $M:=E\cap K$.
Then restriction gives
\[
\Gal(EK/K)\cong\Gal(E/M),
\]
and, as $K$-algebras,
\[
E\otimes_FK\cong\prod_{i=1}^{[M:F]}EK.
\]
In particular, $E\otimes_FK$ is a field if and only if $E$ and $K$ are
linearly disjoint over $F$.
\end{lemma}

\begin{proof}
Write $E=F[\alpha]$ and let $g\in F[x]$ be the minimal polynomial of
$\alpha$. The translation theorem identifies $\Gal(EK/K)$ with the subgroup
of $\Gal(E/F)$ fixing $E\cap K=M$, hence with $\Gal(E/M)$.

Factor $g=g_1\cdots g_m$ into monic irreducible polynomials over $K$.
Adjoining any root of $g$ to $K$ produces $EK$, so every $g_i$ has degree
$[E:M]$. Consequently $m=[M:F]$, and
\[
E\otimes_FK\cong K[x]/(g)
  \cong\prod_{i=1}^{[M:F]}K[x]/(g_i)
  \cong\prod_{i=1}^{[M:F]}EK.
\]
The product is a field exactly when $M=F$.
\end{proof}

Applying the lemma with $K=K_{C'}$ gives
\[
E\cap K_{C'}=F,
\qquad
\Gal(EK_{C'}/K_{C'})\cong\Gal(E/F)=H.
\]

{\color{blue}
For the following valuation-tower analysis, assume additionally that
\[
V_{C'}\cap K_C=V_C.
\]
This is an extra hypothesis, not a consequence of the blowup construction,
as the preceding example shows. Without it, (IH$_{r-1}$) concerns a
different valuation of $K_C$ from the restriction of $V_{C'}$, and cannot
be used in the displayed tower.
{\color{red}Alternatively, one may choose a valuation
of $K_{C'}$ extending $V_C$, but then (AS), which concerns $V_{C'}$, gives
no unramifiedness statement at that chosen valuation without a further
argument.}
The comparison below is therefore conditional on compatibility.
\par}

Let $V_F:=V_{C'}|_F$. Choose a place $v_E$ of $E$ above $V_F$ and a place
$v_\star$ of $EK_{C'}$ above both $v_E$ and $V_{C'}$. The relevant field
and valuation towers are
\[
\begin{tikzcd}[row sep=1.6em, column sep=1.4em]
 &EK_{C'}\arrow[dl,dash]\arrow[dr,dash,"H"]&\\
 E\arrow[dr,dash,"H"']&&K_{C'}\arrow[dl,dash]\\
 &F\arrow[d,dash,"G/H"]&\\
 &K_C&
\end{tikzcd}
\qquad
\begin{tikzcd}[row sep=1.6em, column sep=1.4em]
 &v_\star\arrow[dl,dash]\arrow[dr,dash]&\\
 v_E\arrow[dr,dash]&&V_{C'}\arrow[dl,dash]\\
 &V_F\arrow[d,dash]&\\
 &V_C&.
\end{tikzcd}
\]

{\color{red}
We can now see exactly what each part of the proposed induction contributes:
\par}
{\color{blue}
Under this additional hypothesis, the two proposed inputs would contribute
the following statements:
\par}
\begin{description}
  \item[\textnormal{(IH$_{r-1}$):}]
    $F/K_C$ is unramified at $V_C$. In particular,
    $e(V_F\mid V_C)=1$.
  \item[\textnormal{(AS):}]
    $EK_{C'}/K_{C'}$ is unramified at $V_{C'}$.
  \item[Missing step:]
    one must prove that $E/F$ is unramified at $V_F$.
  \item[Goal:]
    once the missing step holds, the tower
    $K_C\subset F\subset E$ shows that $E/K_C$ is unramified at $V_C$.
\end{description}

{\color{red}
The gap lies between the second and third items. It is measured by the
following inertia comparison.
\par}
{\color{blue}
Even under the compatibility hypothesis, a gap remains between the second
and third items. It is measured by the following inertia comparison.
\par}

\begin{lemma}\label{lem:inertia-injection}
Let $E/F$ be finite Galois and linearly disjoint from $K/F$. Let $w$ be a
place of $EK$ and write $w_K,w_E,w_F$ for its restrictions. Under the
restriction isomorphism $\Gal(EK/K)\cong\Gal(E/F)$ there is an injection
\[
I(w\mid w_K)\hookrightarrow I(w_E\mid w_F).
\]
After completion, the source is the inertia group of the local compositum
$\widehat E\,\widehat K/\widehat K$. In particular, if
$\widehat E\subseteq\widehat K$, then $I(w\mid w_K)=1$, regardless of the
ramification of $E/F$ at $w_F$.
\end{lemma}

\begin{proof}
Let $\sigma\in I(w\mid w_K)$. It fixes $K$, hence $F$, and restriction gives
$\sigma|_E\in\Gal(E/F)$. Because $\sigma$ preserves $w$, its restriction
preserves $w_E$. Moreover, $\sigma$ acts trivially on the residue field at
$w$, and therefore on the subfield given by the residue of $E$. Thus
$\sigma|_E\in I(w_E\mid w_F)$. Injectivity follows from the restriction
isomorphism. The completed statement is the usual identification of the
decomposition and inertia groups with those of the corresponding extension
of completed fields.
\end{proof}

Applied to the tower above, the lemma gives
\[
I(v_\star\mid V_{C'})\hookrightarrow I(v_E\mid V_F).
\]
Statement (AS) makes the group on the left trivial, whereas the missing step
requires the group on the right to be trivial. An injection with trivial
source gives no such conclusion. In other words, the implication
\[
EK_{C'}/K_{C'}\text{ unramified at }V_{C'}
\quad\Longrightarrow\quad
E/F\text{ unramified at }V_F
\]
is not formal, as the following example shows.

\begin{example}[Absorption of an Artin--Schreier layer]
\label{ex:p2-counterexample}
\label{ex:absorption-explicit}
Let $k=\bar{\F}_p$ and put $F=k((t))$. Define $K=k((\pi))$ and embed
$F\hookrightarrow K$ by
\[
t=\frac{\pi^p}{1-\pi^{p-1}}.
\]
Equivalently,
\[
t^{-1}=\pi^{-p}-\pi^{-1}
       =\wp(\pi^{-1}),
\qquad \wp(z)=z^p-z.
\]
The element $t^{-1}$ has reduced pole order one over $F$. Hence the
Artin--Schreier extension
\[
E=F(y),\qquad y^p-y=t^{-1},
\]
is totally wildly ramified of degree $p$ and has break one by
\Cref{theorem:artin_schreier_ramification}. In particular,
\[
I(E/F)=\Z/p\Z.
\]

Over $K$, however, $\pi^{-1}$ is already a solution of the same equation.
The polynomial $X^p-X-t^{-1}$ is irreducible over $F$, because $t^{-1}$
has reduced pole order one, and hence $F(\pi^{-1})/F$ has degree $p$.
On the other hand, the displayed relation between $t$ and $\pi$ shows that
$K/F$ has degree at most $p$. It follows that $F(\pi^{-1})=K$. Therefore
the fields fit into the diagram
\[
\begin{tikzcd}[row sep=2.2em, column sep=4.5em]
E=F(y)\ar[rr,"\sim","{y\mapsto\pi^{-1}}"']
  &&K=F(\pi^{-1})\\
  &F\ar[ul,hookrightarrow,"p"]\ar[ur,hookrightarrow,"p"']&
\end{tikzcd}
\]
After using the horizontal $F$-isomorphism to identify the two upper fields
inside a common algebraic closure of $F$, we may regard $E$ and $K$ as the
same subfield. Consequently
\[
EK/K=K/K
\]
is trivial and therefore unramified. The inertia comparison is exactly
\[
\underbrace{I(EK/K)}_{1}
\hookrightarrow
\underbrace{I(E/F)}_{\Z/p\Z}.
\]
\end{example}

The example shows directly why the missing implication is not a formal
property of base change: $E/F$ is wildly ramified, whereas after the wild
base change $K/F$, the extension $EK/K$ is trivial. We do not claim that
this particular extension $K/F$ arises from the comparison map
$q_{\cP/H}$, so this is not a counterexample to part~(3) of
\Cref{theorem:cyclic_cases_after_blowup}. It shows only that the desired
conclusion requires more than a general base-change argument.
{\color{red}The direct
Witt-vector calculation supplies that additional input by treating the full
$\Z/p^r\Z$-torsor at once.
\par}

{\color{blue}
The proposed induction therefore encounters two separate obstacles. In
general, the generic exceptional point $\eta_{C'}$ does not necessarily
lie above $\eta_C$, so the two intended inputs concern valuations that
do not fit into the required tower. Even if this compatibility is assumed
and both inputs are granted, unramifiedness after the comparison base
change does not formally imply unramifiedness before it. Thus the two
inputs do not by themselves establish the proposed induction step.
\par}
